\documentclass[/work/src/main.tex]{subfiles}

\begin{document}

\section{Conclusión}

\subsection{Resumen del Proyecto}
Este plantilla de LaTeX demuestra cómo estructurar un documento profesional utilizando herramientas modernas:

\begin{itemize}
    \item \textbf{Dev Containers:} Ambiente reproducible con Docker
    \item \textbf{Subfiles:} Organización modular del código
    \item \textbf{LuaLaTeX:} Motor de compilación moderno
    \item \textbf{BibLaTeX + Biber:} Gestión bibliográfica avanzada
    \item \textbf{VS Code:} Editor con LaTeX Workshop integrado
\end{itemize}

\subsection{Ventajas del Enfoque}
La estructura implementada ofrece varios beneficios:

\subsubsection{Modularidad}
Cada sección es independiente y puede compilarse por separado. Esto acelera el desarrollo y facilita el trabajo en colaboración.

\subsubsection{Reproducibilidad}
El uso de Dev Containers garantiza que cualquier persona pueda replicar exactamente el mismo entorno de compilación, eliminando problemas de ``funciona en mi máquina''.

\subsubsection{Profesionalismo}
La estructura clara y la separación de concerns (configuración, contenido, ejemplos) hacen el proyecto fácil de mantener y extender.

\subsubsection{Automatización}
Con latexmk configurado correctamente, la compilación es completamente automática y genera PDFs limpios con auxiliares organizados.

\subsection{Próximas Mejoras}
Algunas extensiones posibles para este proyecto:

\begin{enumerate}
    \item Añadir más secciones temáticas
    \item Crear un sistema de comentarios internos (paquete \texttt{todonotes})
    \item Implementar control de versiones con git
    \item Publicación automática en GitHub / Overleaf
    \item Generación de documentación desde el código
\end{enumerate}

\subsection{Agradecimientos}
Este proyecto fue construido utilizando tecnologías y paquetes de código abierto. Agradecemos especialmente a:

\begin{itemize}
    \item \textbf{The LaTeX Project} — por la excelente herramienta de composición
    \item \textbf{James Yu} — por el paquete LaTeX Workshop
    \item \textbf{Docker} — por los Dev Containers
    \item \textbf{La comunidad TeX/LaTeX} — por el soporte y documentación
\end{itemize}

\end{document}
